% ============================================================================
% PART II: PHYSICAL APPLICATIONS
% Chapter 13: Topological Quantum Computation
% ============================================================================
%
% Purpose
%   Non-abelian anyons as the substrate for fault-tolerant quantum
%   computation. Kitaev's proposal, Ising worked example, honesty box.
%
% Status: NEW WRITING. Guided by BSIEW tasks 13_1--13_3.
%
% Length target: ~8 pages
%
% Section plan
%   13.1  Non-abelian anyons and Kitaev's proposal
%   13.2  Ising anyons: four sigmas = two qubits
%   13.3  Honesty box: experiment vs. theory gap
%   13.4  Brief outlook: what a working platform would mean
%
% Owner: Helena
% Stage: V
%
% BSIEW task files:
%   13_1_non_abelian_kitaev.md
%   13_2_ising_gates.md
%   13_3_honesty_box.md
%
% References
%   Kitaev 2003 (\cite{Kitaev2003FaultTolerant})
%   Nayak et al. 2008 (\cite{NayakSimonSternFreedmanDasSarma2008NonabelianAnyons})
%
% Dependencies
%   - Ch 10 (anyon formalism, Ising data from comparison table).
%   - Ch 12 (Andersen 2023 context).

\section{Topological Quantum Computation}
\label{sec:tqc}

\subsection{Non-Abelian Anyons and Kitaev's Proposal}
\label{subsec:kitaev-proposal}

% TODO (Task 13_1): ~1.5 pages.
% 1. Topological degeneracy: N non-abelian anyons of fixed fusion type have
%    degenerate ground space, dimension grows exponentially in N.
% 2. Topological protection: local perturbations commute with ground-space
%    projector to exponential accuracy in system size; logical qubits
%    inaccessible to local noise.
% 3. Braiding as unitary gate: exchanging anyons effects unitary on ground
%    space; computable sequence of braids implements a gate.
% Cite: Kitaev 2003, Nayak et al. RMP 2008.

\subsection{Ising Anyons: Four $\sigma$'s $=$ Two Qubits}
\label{subsec:ising-gates}

% TODO (Task 13_2): ~3 pages. WORKED EXAMPLE.
% 1. Ising fusion: sigma x sigma = 1 + psi, so 4 sigma's fuse to
%    4-dim space = 2 qubits.
% 2. Braiding group: exchanges of adjacent sigma's generate representation
%    of braid group B_4 on 2-qubit space.
% 3. Gates implemented: generated gate set is Clifford group --- powerful
%    but not universal.
% 4. Need for additional resources: topologically-protected Clifford +
%    non-topological T gate gives universality. Fault-tolerance comes
%    from magic-state distillation.
% Cite: Nayak et al. RMP 2008.

\subsection{Required Honesty Box}
\label{subsec:tqc-honesty-box}

% TODO (Task 13_3): ~1 page. REQUIRED BY RUBRIC.
%
% Honesty box (visually distinct: tcolorbox or similar) stating:
%
% 1. No natural non-abelian anyonic phase has been unambiguously confirmed.
%    nu = 5/2 Moore-Read is strongest candidate but not definitive.
% 2. Majorana-based proposals: semiconductor/superconductor hybrids have
%    produced suggestive but contested signatures; prominent claims retracted.
%    Field has lower confidence than 5 years ago.
% 3. Synthetic realizations (Ch 12 Sec 12.3) implement braid algebra but
%    not passive topological protection.
% 4. Fault-tolerance threshold: even with a working platform, achieving
%    threshold requires substantial improvements in coherence times and
%    operational fidelity.
%
% Tone: matter-of-fact, not hedging. The honesty is the point.

\subsection{Brief Outlook: What a Working Platform Would Mean}
\label{subsec:tqc-outlook}

% TODO (new subsection): ~1 page.
% Connect back to Ch 10 anyon data: which systems (Ising, Fibonacci) would
% give universal TQC, and what the experimental requirements are.
% Ising: Clifford only, needs magic state distillation.
% Fibonacci: universal from braiding alone, but no known natural realization.
% Cite: Nayak et al. RMP 2008.
